% APÊNDICES -- VERSÃO CONCISA TCC UESC
% ==========================================================================
\begin{apendicesenv}

\chapter{Ficha técnica e especificação dimensional do modelo de voo qualificado (Projeto \#28)}
\label{apendice:ficha-tecnica}

Este apêndice consolida as especificações geométricas, mecânicas e normativas do chassi primário CubeSat 1U qualificado para lançamento espacial, sintetizado pelo algoritmo evolutivo NSGA-II e selecionado pelo método de decisão multicritério TOPSIS.

\section*{Especificações geométricas e construtivas (L-PBF)}
\begin{itemize}
  \item \textbf{Fator de Forma e Envelope:} Padrão CalPoly CubeSat 1U ($\SI{100.0}{\milli\metre} \times \SI{100.0}{\milli\metre} \times \SI{113.5}{\milli\metre}$).
  \item \textbf{Trilhos Guia de Implantação (\textit{Deployer Rails}):} Quatro trilhos longitudinais maciços de seção quadrada (\rail{}), com rugosidade $Ra \le \SI{1.6}{\micro\metre}$ e anodização dura.
  \item \textbf{Ângulo de Reentrância Celular ($\theta$):} $\SI{-21.40}{\degree}$ ($\theta_{\text{overhang}} = \SI{68.60}{\degree} \ge \SI{35.0}{\degree}$ perante a mesa de construção, autossuportado sem suportes internos).
  \item \textbf{Espessura de Costela Celular ($t$):} $\SI{0.613}{\milli\metre}$ ($> \SI{0.50}{\milli\metre}$, consolidada por feixe laser de alta resolução).
  \item \textbf{Comprimento da Haste Oblíqua ($l$):} $l^* = \SI{10.42}{\milli\metre}$ entre eixos nodais ($l_{\text{free}} = \SI{6.42}{\milli\metre}$ de segmento livre de flexão).
  \item \textbf{Altura da Parede Vertical ($h$):} $h^* = \SI{14.85}{\milli\metre}$ entre eixos nodais ($h_{\text{free}} = \SI{8.75}{\milli\metre}$ de segmento livre).
  \item \textbf{Raio de Concordância nos Entalhes ($R_{\text{notch}}$):} $\SI{0.50}{\milli\metre}$ ($K_t \approx \num{1.26}$).
  \item \textbf{Folga Livre de Despoeiramento:} $\text{gap}_{\text{despoeiramento}} = 2\,l\cos\theta - 2\,t = \SI{10.74}{\milli\metre}$ (passagem livre local $\ge \SI{7.40}{\milli\metre}$, superando o teto de $\SI{1.50}{\milli\metre}$).
  \item \textbf{Canais de Drenagem de Pó Residual:} Orifícios circulares desobstruídos com diâmetro $d_{\text{drain}} \ge \SI{2.00}{\milli\metre}$.
\end{itemize}

\section*{Propriedades constitutivas da liga aeroespacial AlSi10Mg (L-PBF)}
\begin{itemize}
  \item \textbf{Tratamento Térmico Pós-Processamento:} Rota canônica a $\SI{300}{\celsius} \pm \SI{10}{\celsius}$ por $\SI{2.0}{\hour}$ ($\sigma_{\text{yield}} \ge \SI{230.0}{\mega\pascal}$) ou rota mitigadora a $\SI{250}{\celsius} / \SI{2.0}{\hour}$ ($\sigma_{\text{yield}} \ge \SI{270.0}{\mega\pascal}$), com pré-aquecimento de mesa a $T_{\text{bed}} = \SI{150}{\celsius}$.
  \item \textbf{Tolerância a Defeitos e Fratura (Murakami LEFM):} $\Delta K_{\text{defect}} = \SI{1.28}{\mega\pascal\sqrt{\metre}} < \Delta K_{\text{th}} = \SI{2.50}{\mega\pascal\sqrt{\metre}}$ para porosidades típicas $d_{\text{pore}} \le \SI{50.0}{\micro\metre}$ ($a_{\text{th}} \approx \SI{169.0}{\micro\metre}$).
  \item \textbf{Massa Específica ($\rho$):} $\SI{2670.0}{\kilo\gram\per\cubic\metre}$.
  \item \textbf{Módulo de Elasticidade Longitudinal ($E$):} $\SI{68.0}{\giga\pascal}$.
  \item \textbf{Coeficiente de Poisson do Metal Sólido ($\nu$):} $\num{0.33}$.
  \item \textbf{Limite de Escoamento a $0{,}2\%$ ($\sigma_{\text{yield}}$):} $\SI{230.0}{\mega\pascal}$.
  \item \textbf{Limite de Resistência à Tração ($\sigma_{\text{ult}}$):} $\SI{340.0}{\mega\pascal}$.
  \item \textbf{Densidade Relativa do Núcleo Celular ($\rho_{\text{rel}}$):} $\num{0.148}$ ($\SI{14.80}{\percent}$).
  \item \textbf{Coeficiente de Poisson Efetivo Aparente ($\nu_{\text{eff}}$):} $\num{-0.428}$ (comportamento auxético bidirecional).
\end{itemize}

\section*{Indicadores de desempenho e qualificação dinâmica (NASA GEVS)}
\begin{itemize}
  \item \textbf{Massa Estrutural Total:} $\SI{0.11554}{\kilo\gram}$ ($\approx \SI{0.116}{\kilo\gram}$), representando economia de $\SI{62.49}{\percent}$ frente ao bloco maciço ($\SI{0.308}{\kilo\gram}$).
  \item \textbf{Primeira Frequência Natural ($f_1$, engastado / \code{'clamped'}):} $\SI{579.62}{\hertz}$ ($\SI{+13.20}{\percent}$ frente aos $\SI{512.0}{\hertz}$ do maciço).
  \item \textbf{Primeira Frequência Natural ($f_1$, suporte flexível / \code{'test_pod'}):} $\SI{304.10}{\hertz}$ ($MS_{f_1} = +2{,}04$ frente a $\SI{100.0}{\hertz}$ da GEVS).
  \item \textbf{Transmissibilidade Dinâmica Passiva ($T$):} $\num{0.1997}$ ($\SI{76.50}{\percent}$ ou $\SI{12.6}{\decibel}$ de atenuação na carga útil).
  \item \textbf{Aceleração Eficaz na Carga Útil ($\Grmspay$):} $\SI{2.82}{\Grms}$ (sob excitação de base de $\SI{14.10}{\Grms}$).
  \item \textbf{Tensão Equivalente de Pico a $3\sigma$ ($\sigma_{3\sigma}$):} $\SI{166.91}{\mega\pascal}$.
  \item \textbf{Margem de Segurança de Escoamento ($MS_{\text{yield}}$):} $\num{+0.102} > 0{,}0$ ($FS_{\text{yield}} = \num{1.25}$, $\sigma_{\text{adm, yield}} = \SI{184.0}{\mega\pascal}$, NASA-STD-5020B).
  \item \textbf{Margem de Segurança de Ruptura ($MS_{\text{ult}}$):} $\num{+0.455} > 0{,}0$ ($FS_{\text{ult}} = \num{1.40}$, $\sigma_{\text{adm, ult}} = \SI{242.86}{\mega\pascal}$).
  \item \textbf{Dano Acumulado de Fadiga de Steinberg ($D$):} $\num{0.0458} \ll \num{0.25}$ (fator de segurança de vida de $5{,}4\times$, NASA-HDBK-7005).
  \item \textbf{Incerteza de Malha pelo GCI de Roache (ASME V\&V 20):} $GCI_{\text{fine}}^{f_1} = \SI{0.23}{\percent} \ll \SI{1.5}{\percent}$ (Tet10 de alta ordem com $\num{342800}$ nós).
\end{itemize}

\end{apendicesenv}
